\documentclass[11pt,a4paper]{article}
\usepackage[T1]{fontenc}
\usepackage{lmodern,amsmath,amssymb,amsthm,mathtools,microtype}
\usepackage[margin=25mm]{geometry}
\usepackage{xcolor}
\usepackage[colorlinks=true,linkcolor=blue!45!black,citecolor=blue!45!black,urlcolor=blue!45!black]{hyperref}
\newtheorem{theorem}{Theorem}[section]
\newtheorem{lemma}[theorem]{Lemma}
\newtheorem{proposition}[theorem]{Proposition}
\setlength{\emergencystretch}{2em}
\title{Explicit Quadratic Height Bounds\\for Elliptic Multiples}
\author{}\date{8 November 2026}
\begin{document}\maketitle
\begin{abstract}
For a specified point on an elliptic curve over $\mathbb Q(\sqrt2)$,
we prove an explicit upper height bound quadratic in its integer index.
The proof uses exact duplication and differential-addition identities,
an elementary induction, and a comparison between absolute height and
primitive coefficient height. We also bound the coefficients of finite
products excluding these elliptic abscissas. The results are
unconditional and require no determination of the full Mordell--Weil group.
\end{abstract}
\section{Quadratic height bounds for the fixed elliptic multiples}
\label{sec:quadratic-fixed-height}

Put $K=\mathbb Q(\sqrt2)$ and
\[
 D=11-6\sqrt2,\qquad E:\ y^2=x^3-Dx,\qquad
 P=\left(\frac{57-28\sqrt2}{8},
          \frac{-660+417\sqrt2}{32}\right).
\]
The height $h$ below is the absolute logarithmic Weil height. If
$x=(a+b\sqrt2)/c$ with $a,b,c\in\mathbb Z$, $c>0$, and
$\gcd(a,b,c)=1$, put
\[
 h_0(x)=\left\lceil\log_2\max(1,|a|,|b|,c)\right\rceil.
\]

\begin{lemma}\label{lem:fixed-height-nontorsion}
The point $P$ belongs to $E(K)$ and has infinite order.
\end{lemma}
\begin{proof}
Substitution verifies the curve equation. At the unique place above two,
normalize $v$ by $v(2)=1$. The displayed coordinates have valuations
$v(x(P))=-3$ and $v(y(P))=-9/2$, so the formal parameter
$z=-x/y$ has $v(z(P))=3/2$. The integral formal group of this
Weierstrass equation is odd: for every integer $m$,
$[m](Z)=mZ+Z^3g_m(Z^2)$ with integral coefficients.
At a parameter of valuation at least $3/2$, multiplication by an odd
integer preserves its valuation, and multiplication by two raises it
by one. In each case the leading term has strictly smaller valuation
than the remaining terms. Thus
\[
                   v(z([n]P))=3/2+v_2(n)
                   \qquad(n\ne0).
\]
It is finite, and hence $[n]P\ne\mathcal O$ for every nonzero $n$.
In particular all its abscissas are finite and nonzero.
\end{proof}

\begin{lemma}\label{lem:fixed-height-recurrences}
For $n\ge1$, write $h_n=h(x([n]P))$. Then
\[
 \begin{aligned}
 h_{2n}&\le4h_n+\log392,\\
 h_{2n+1}&\le2h_n+2h_{n+1}+\log8036.
 \end{aligned}
\]
The second formula is used for $n\ge1$.
\end{lemma}
\begin{proof}
The duplication and differential-addition identities are
\[
 x(2Q)=\frac{(x(Q)^2+D)^2}{4x(Q)(x(Q)^2-D)},
 \qquad
 x(Q+R)x(Q-R)=
       \left(\frac{x(Q)x(R)+D}{x(Q)-x(R)}\right)^2.
\]
For the second identity, the usual addition formulas reduce it to
\[
 (u+v)^2(uv-D)^2-4uv(u^2-D)(v^2-D)
                    =(u-v)^2(uv+D)^2,
\]
which is a polynomial identity. Taking $Q=[n]P$ and
$R=[n+1]P$ makes $x(Q-R)=x(P)$. Its denominator is nonzero:
an equality $x([n]P)=x([n+1]P)$ would contradict the infinite
order of $P$.

For each place, bound a homogeneous polynomial by the largest
coordinate norm to its degree. At a finite place the coefficient
bound is the maximum coefficient norm; at an infinite place use
the sum of the coefficient norms. Applied to duplication, whose
homogeneous numerator and denominator have degree four, this gives
\[
             h(x(2Q))\le4h(x(Q))+2h(D)+\log8.
\]
Applied to differential addition, with bidegree $(2,2)$ in the two
abscissas and denominator $x(P)(x(Q)-x(R))^2$, it gives
\[
 h(x(Q+R))\le2h(x(Q))+2h(x(R))+2h(D)+h(x(P))+\log4.
\]
For example, the archimedean coefficient sums for the second
formula are bounded by
$4\max(1,|D|)^2\max(1,|x(P)|)$, and the corresponding finite-place
bound has no factor four. These local estimates sum to the stated
absolute-height estimates.

Both real conjugates of $D$ exceed one, and $N_{K/\mathbb Q}(D)=49$,
so $h(D)=\log7$. Both real conjugates of $x(P)$ exceed one, its
norm is $41^2/8^2$, and its denominator ideal has norm $8^2$.
Consequently $h(x(P))=\log41$. Substitution gives
$2h(D)+\log8=\log392$ and
$2h(D)+h(x(P))+\log4=\log8036$.
\end{proof}

\begin{theorem}\label{thm:fixed-height-quadratic}
For every integer $n\ge1$,
\[
 h(x([n]P))\le n^2\log23534-\log574,
 \qquad h_0(x([n]P))\le30n^2.
\]
These inequalities are unconditional for the displayed curve and point.
\end{theorem}
\begin{proof}
Write $A=\log23534$ and $B=\log574$. The integer identities
\[
 23534=41\cdot574,\qquad
 574^3=23534\cdot8036,\qquad574^3\ge392
\]
give $h_1=A-B$, $3B=A+\log8036$, and $3B\ge\log392$.
Induction using Lemma~\ref{lem:fixed-height-recurrences} proves
$h_n\le An^2-B$. For an even index,
\[
 h_{2n}\le4(An^2-B)+\log392\le A(2n)^2-B.
\]
For an odd index at least three,
\[
 \begin{split}
 h_{2n+1}&\le2(An^2-B)+2(A(n+1)^2-B)+\log8036\\
 &=A(2n+1)^2+A-4B+\log8036
  =A(2n+1)^2-B.
 \end{split}
\]

We next compare coefficient height with absolute height. Write
$x=(a+b\sqrt2)/c$ in the primitive form above, and let
$\mathfrak d=\{u\in\mathcal O_K:ux\in\mathcal O_K\}$ be its
denominator ideal. Since $\mathcal O_K=\mathbb Z[\sqrt2]$, the
integer $c$ is the least positive rational integer with
$cx\in\mathcal O_K$. At a prime above a rational prime $p$, use the integer-valued ideal
valuation $\operatorname{ord}_{\mathfrak p}$, so that
$\operatorname{ord}_{\mathfrak p}(p)=e_{\mathfrak p}$, and put
$d_{\mathfrak p}=\max(0,-\operatorname{ord}_{\mathfrak p}(x))$.
The $p$-adic exponent of $N\mathfrak d$ is
$\sum_{\mathfrak p\mid p}f_{\mathfrak p}d_{\mathfrak p}$.
Multiplying by a ramification index $e_{\mathfrak p}$ makes this
at least $d_{\mathfrak p}$. Therefore
$(N\mathfrak d)x$ is integral and $c\mid N\mathfrak d$.
If $h_{\mathrm{fin}}(x)$ denotes the finite-place part of $h(x)$,
then $N\mathfrak d=\exp(2h_{\mathrm{fin}}(x))$.

Writing $x'$ for the real conjugate, one has
\[
 |a|=c\left|\frac{x+x'}2\right|\le c\max(|x|,|x'|),\qquad
 |b|=c\left|\frac{x-x'}{2\sqrt2}\right|
                         \le c\max(|x|,|x'|).
\]
The product $\max(1,|x|)\max(1,|x'|)$ equals
$\exp(2h_{\infty}(x))$. It follows that
\[
                  \max(1,|a|,|b|,c)\le\exp(2h(x)).
\]
Finally $23534<2^{15}$, so the first height bound gives
$\exp(2h_n)<2^{30n^2}$ and hence $h_0(x([n]P))\le30n^2$.
\end{proof}

\subsection{The indexed exclusion polynomial}

Let $\mathcal A$ be a finite free algebra over an integer polynomial
ring, with a fixed basis. For a multiplication matrix use the maximum
column sum of the coefficient $\ell^1$-norms. This norm is
submultiplicative. Suppose $b\ge0$ and the multiplication matrices of
$N_0,N_1,D_0\in\mathcal A$ have norm at most $2^b$.
Write $x([n]P)=(a_n+b_n\sqrt2)/c_n$ in primitive form and set
\[
 \mathcal E_n=(c_nN_0-a_nD_0)^2-2(c_nN_1-b_nD_0)^2,
 \qquad \mathcal E_{\le M}=\prod_{n=1}^M\mathcal E_n.
\]
Each fixed triple is computed by exact elliptic arithmetic. The product
is finite and is a literal polynomial specification, independently of
whether its coefficients are expanded.

\begin{proposition}\label{prop:fixed-height-cubic-product}
For every integer $M\ge1$,
\[
 \|M_{\mathcal E_{\le M}}\|
       \le 2^{20M^3+30M^2+(2b+14)M}
       \le 2^{(2b+64)M^3}.
\]
If $b\le1048$ and $M\ge64$, the first exponent is at most $21M^3$.
\end{proposition}
\begin{proof}
Theorem~\ref{thm:fixed-height-quadratic} and the triangle inequality give
\[
 \|M_{c_nN_0-a_nD_0}\|,\ \|M_{c_nN_1-b_nD_0}\|
       \le2^{30n^2+b+1}.
\]
Squaring and summing the two terms of $\mathcal E_n$ therefore
gives $\|M_{\mathcal E_n}\|\le2^{60n^2+2b+4}$.
Sum these logarithmic bounds and use
$\sum_{n=1}^M n^2=M(M+1)(2M+1)/6$.
The second bound follows from $M,M^2\le M^3$.
For the last assertion, $2b+14\le2110$, and
\[
                   30\cdot64+2110=4030<64^2.
\]
The function $M^2-30M$ is increasing for integers $M\ge64$.
Thus $30M^2+(2b+14)M\le M^3$, proving the result.
\end{proof}

Under an algebra homomorphism $\mathcal A\to\mathbb Q$ with
the image of $D_0$ nonzero, the factor
$\mathcal E_n$ vanishes exactly when
$(N_0+N_1\sqrt2)/D_0=x([n]P)$. Indeed a rational pair satisfies
$u^2-2v^2=0$ only when both entries vanish. The displayed product therefore gives an exact polynomial exclusion
of these finitely many abscissas, with the proved coefficient bound.


\end{document}
